\section*{Capítulo \ref{ch:b3:many-electron-atoms} --- Átomos polielectrónicos y símbolos de término}

\begin{solution}{exo:b3:many-electron-atoms:1}
\[
\Psi = \frac{1}{\sqrt6}\begin{vmatrix}1s\alpha(1) & 1s\beta(1) & 2s\alpha(1)\\
1s\alpha(2) & 1s\beta(2) & 2s\alpha(2)\\ 1s\alpha(3) & 1s\beta(3) & 2s\alpha(3)
\end{vmatrix}.
\]
$1s$ solo ofrece dos \omterm{def:b3:many-electron-atoms:spin-orbital}{espín-orbitales}, $1s\alpha$ y $1s\beta$; un tercer electrón
repetiría uno de ellos, con lo que dos columnas serían iguales: el determinante es nulo.
\end{solution}

\begin{solution}{exo:b3:many-electron-atoms:2}
$\binom62 = 15$, $\binom63 = 20$, $\binom{10}{2} = 45$. $p^3$: ${}^4$S (4) +
${}^2$D (10) + ${}^2$P (6) = 20. $d^2$: ${}^3$F (21) + ${}^3$P (9) + ${}^1$G (9) +
${}^1$D (5) + ${}^1$S (1) = 45.
\end{solution}

\begin{solution}{exo:b3:many-electron-atoms:3}
N ($2p^3$, semillena): \termsym{4}{S}{3/2}. O ($2p^4$, más que semillena):
\termsym{3}{P}{2}. F ($2p^5$): \termsym{2}{P}{3/2}. \ce{Ti^2+} ($3d^2$):
\termsym{3}{F}{2}. \ce{Fe^3+} ($3d^5$, semillena, todos los espines paralelos, $L = 0$):
\termsym{6}{S}{5/2}.
\end{solution}

\begin{solution}{exo:b3:many-electron-atoms:4}
${}^2$D: $J = \frac52$ (6 estados) y $\frac32$ (4); $6 + 4 = 10 = 5 \times 2$.
${}^3$P: $J = 2$ (5), 1 (3), 0 (1); $5 + 3 + 1 = 9 = 3 \times 3$.
\end{solution}

\begin{solution}{exo:b3:many-electron-atoms:5}
Intervalos de $16.4$ ($J = 1 - 0$) y $27.0$ ($2 - 1$): razón 1.65 en lugar de 2.
$A = 16.4/1 = \qty{16.4}{cm^{-1}}$ a partir del primero, $27.0/2 = \qty{13.5}{cm^{-1}}$
a partir del segundo. Un único $A$ no basta: los niveles están ligeramente mezclados con
los de otros términos (${}^1$D, ${}^1$S), y el \omterm{def:b3:many-electron-atoms:russell-saunders}{acoplamiento de Russell--Saunders} es una
aproximación.
\end{solution}

\begin{solution}{exo:b3:many-electron-atoms:6}
$\lambda = 10^7/\tilde\nu$: \qty{589.76}{nm} (D$_1$) y \qty{589.16}{nm}
(D$_2$). Separación: $\qty{17.20}{cm^{-1}} = \qty{2.13}{meV}$. $E(\frac32) -
E(\frac12) = \frac32A$, luego $A = \qty{11.47}{cm^{-1}}$.
\end{solution}

\begin{solution}{exo:b3:many-electron-atoms:7}
$3p \to 3s$: permitida ($\Delta l = -1$). $3d \to 3s$: prohibida ($\Delta l = -2$).
$3d \to 3p$: permitida. $4s \to 3s$: prohibida ($\Delta l = 0$, sin cambio de
paridad). $4s \to 3p$: permitida.
\end{solution}

\begin{solution}{exo:b3:many-electron-atoms:8}
Media de ${}^3$P: $(5 \times 169086.76 + 3 \times 169086.84 + 169087.83)/9 =
\qty{169086.91}{cm^{-1}}$. Separación hasta ${}^1$P: \qty{2047.98}{cm^{-1}} $=
\qty{0.254}{eV}$, luego $K(1s,2p) = \qty{0.127}{eV}$, tres veces menor que
$K(1s,2s)$. La \omterm{def:b3:many-electron-atoms:exchange}{integral de intercambio} mide el solapamiento de las densidades de los dos
orbitales; el orbital $2s$ penetra en la región $1s$ (tiene densidad
cerca del núcleo), mientras que el orbital $2p$ se anula en el núcleo y solapa menos.
\end{solution}

\begin{solution}{exo:b3:many-electron-atoms:9}
El mayor valor, $M_L = 4$, exige los dos electrones en $m_l = 2$, con espines opuestos: un
singlete, ${}^1$G (9 estados). El siguiente, $M_L = 3$, procede de $m_l = 2, 1$, con
$M_S = 1$ posible: ${}^3$F (21). Queda $M_L = 2$ solo con $M_S = 0$:
${}^1$D (5). Queda $M_L = 1$ con $M_S = 1$: ${}^3$P (9). Queda un estado con
$M_L = M_S = 0$: ${}^1$S. Total: 45.
\end{solution}

\begin{solution}{exo:b3:many-electron-atoms:10}
En $\mathbf r_1 = \mathbf r_2 = \mathbf r$, la función triplete es
$\frac{1}{\sqrt2}[1s(\mathbf r)2s(\mathbf r) - 2s(\mathbf r)1s(\mathbf r)] = 0$; la
función singlete es $\sqrt2\,1s(\mathbf r)2s(\mathbf r) \ne 0$. En el triplete, los
electrones nunca se encuentran en el mismo punto y están por término medio más
alejados, de modo que su repulsión es menor: $E_+ - E_- = 2K > 0$.
\end{solution}

\begin{solution}{exo:b3:many-electron-atoms:11}
$d^8$ está más que semillena: el acoplamiento es invertido, y $J = L + S = 4$ es
el más bajo. Intervalos: $E(3) - E(4) = 1360.7$ y $E(2) - E(3) = 908.9$; razón
1.50, frente a los $4/3 = 1.33$ de Landé. $|A| = 1360.7/4 = \qty{340}{cm^{-1}}$
($A < 0$). El \omterm{def:b3:many-electron-atoms:spin-orbit}{acoplamiento espín--órbita} es mucho mayor en este ion más pesado que en el
carbono (\qty{16}{cm^{-1}}).
\end{solution}

\begin{solution}{exo:b3:many-electron-atoms:12}
$\sum_J(2J+1)[J(J+1) - L(L+1) - S(S+1)] = 0$ (los $(2L+1)(2S+1)$ estados pueden
contarse en cualquiera de las dos bases, y $\sum M_LM_S = 0$ sobre los no acoplados).
Para ${}^3$P: $1(0 - 4) + 3(2 - 4) + 5(6 - 4) = -4 - 6 + 10 = 0$. La media
ponderada del carbono es $(0 + 3 \times 16.42 + 5 \times 43.41)/9 = \qty{29.6}{cm^{-1}}$:
la energía que tendría el término sin \omterm{def:b3:many-electron-atoms:spin-orbit}{acoplamiento espín--órbita}.
\end{solution}

\begin{solution}{pb:b3:many-electron-atoms:1}
\textbf{1.} $1s\alpha$, $1s\beta$, $2s\alpha$, $2s\beta$.
\textbf{2.} $|1s\alpha\,2s\alpha|$, $|1s\alpha\,2s\beta|$, $|1s\beta\,2s\alpha|$,
$|1s\beta\,2s\beta|$.
\textbf{3.} $|1s\alpha\,2s\alpha| = \frac{1}{\sqrt2}[1s(1)2s(2) - 2s(1)1s(2)]\,
\alpha(1)\alpha(2)$, desarrollando el determinante $2\times2$.
\textbf{4.} Su suma es $\frac{1}{\sqrt2}[1s(1)2s(2) - 2s(1)1s(2)]\cdot
\frac{1}{\sqrt2}[\alpha(1)\beta(2) + \beta(1)\alpha(2)]$ (multiplicada por $\sqrt2$ y
después normalizada); su diferencia da $\frac{1}{\sqrt2}[1s(1)2s(2) +
2s(1)1s(2)]\cdot\frac{1}{\sqrt2}[\alpha(1)\beta(2) - \beta(1)\alpha(2)]$.
\textbf{5.} Parte espacial simétrica con la función de espín antisimétrica ($S = 0$);
parte espacial antisimétrica con las tres funciones de espín simétricas ($S = 1$).
\textbf{6.} $2S + 1 = 1$ y 3: uno y tres estados de la misma energía.
\textbf{7.} $[\ce{Ar}]\,3d^2$; $\binom{10}{2} = 45$ determinantes.
\textbf{8.} Valores máximos: $S = 1$, $L = 2 + 1 = 3$; menos que semillena:
\termsym{3}{F}{2}, el nivel en 0 de los datos.
\textbf{9.} $21 + 9 + 9 + 5 + 1 = 45$.
\textbf{10.} Normal ($J = 2$ el más bajo). Intervalos de 184.9 y 235.5: razón 1.27
frente a $4/3 = 1.33$; la regla se cumple con un margen del 5\,\%.
\textbf{11.} ${}^3$F: $(5 \times 0 + 7 \times 184.9 + 9 \times 420.4)/21 =
\qty{241.8}{cm^{-1}}$; ${}^3$P: $(10538.4 + 3 \times 10603.6 + 5 \times
10721.2)/9 = \qty{10661.7}{cm^{-1}}$.
\textbf{12.} $15B = \qty{10419.9}{cm^{-1}}$, $B = \qty{695}{cm^{-1}}$.
\textbf{13.} $\Delta S = 0$.
\textbf{14.} $\Delta S = 1$ está prohibida, y $1s2s \to 1s^2$ tiene $\Delta l = 0$
(sin cambio de paridad; ${}^3$S$_1 \to {}^1$S$_0$ tiene además $\Delta L = 0$ entre dos estados
S). El estado es metaestable: vive mucho más que los estados excitados
ordinarios.
\textbf{15.} $169086.91 - 159855.97 = \qty{9230.94}{cm^{-1}}$: \qty{1083.3}{nm}, en
el infrarrojo cercano.
\textbf{16.} $171134.89 - 166277.44 = \qty{4857.45}{cm^{-1}}$: \qty{2058.7}{nm}.
\textbf{17.} $\qty{171134.89}{cm^{-1}}$: \qty{58.43}{nm}, en el ultravioleta de
vacío (absorbido por el aire).
\textbf{18.} Cada familia tiene sus propias rayas y su propio estado más bajo, y las
dos nunca se conectan por la luz: se comportan como dos gases con dos espectros.
\textbf{19.} $E_\pm = E_{1s} + E_{2s} + J \pm K$ (singlete $+$, triplete $-$).
\textbf{20.} El triplete: su función espacial se anula cuando los electrones coinciden,
de modo que se repelen menos.
\textbf{21.} \qty{6421.47}{cm^{-1}}, $6421.47/8065.54 = \qty{0.796}{eV}$.
\textbf{22.} Separación: $\qty{2047.98}{cm^{-1}} = \qty{0.254}{eV}$, luego $K(1s,2p) =
\qty{0.127}{eV}$.
\textbf{23.} El $2s$ penetra en la región $1s$ y la solapa más que el $2p$.
\textbf{24.} Sí: en ambas configuraciones, el triplete (mayor $S$) es el más bajo.
\textbf{25.} $K = \frac12 \times \qty{0.796}{eV}$: \textbf{la \omterm{def:b3:many-electron-atoms:exchange}{integral de
intercambio} $K(1s,2s)$ del helio vale \qty{0.398}{eV}}, medida como la mitad de la
separación singlete--triplete.
\end{solution}
